Text-to-image generative models may reproduce representational patterns present in their training data. This work develops an applied, quantitative, and exploratory analysis of 64 synthetic images depicting university environments, generated by four systems from neutral and inclusive prompts. BiasAuditFW separates processing into two observational units. DeepFace with RetinaFace produces predictions for each detected face, whereas CLIP estimates, once per image, perceived racial diversity and gender-presentation diversity margins. The CLIP scores compare textual prototypes of diversity and homogeneity using explicit L2 normalization. Face detection will be evaluated against bounding boxes annotated by two evaluators. Attribute predictions will also be compared with FairFace as a secondary reference, without treating either classifier as a source of demographic identity. Mann--Whitney and Wilcoxon tests will be applied separately to both margins at image level, with Holm correction. The Patchwork Effect will be discussed only if human ratings, semantic scores, and facial composition provide converging evidence.

\noindent\textbf{Keywords:} bias analysis; generative AI; computer vision; CLIP; perceived diversity; demographic representation.
